% cudaDMA Illegal Instruction Issue Summary
% Include this file in your main document with % cudaDMA Illegal Instruction Issue Summary
% Include this file in your main document with % cudaDMA Illegal Instruction Issue Summary
% Include this file in your main document with % cudaDMA Illegal Instruction Issue Summary
% Include this file in your main document with \input{cudadma-illegal-instruction-issue}

\subsection{cudaDMA Illegal Instruction Issue}

\subsubsection{Problem Description}

The \texttt{jacobi2D\_cudaDMA} kernel encounters an illegal instruction error when compiled with optimizations (\texttt{-O2}/\texttt{-O3}), but runs without errors when compiled with debug flags (\texttt{-g -G}). However, a critical observation reveals suspicious behavior even in the debug build:

\begin{itemize}
    \item \textbf{No illegal instruction error} in debug mode
    \item \textbf{No CPU-GPU mismatch} reported (results appear correct)
    \item \textbf{Print statements inside kernel do NOT appear} in console output
\end{itemize}

This suggests the kernel may not be executing at all, yet falsely reports success.

\subsubsection{Root Causes}

\begin{enumerate}
    \item \textbf{Architecture mismatch:} The \texttt{cudaDMAStrided} helper likely emits SM80+ instructions (e.g., \texttt{cp.async}) that are unsupported on older GPUs. Optimized builds inline these instructions without runtime guards.
    
    \item \textbf{Silent kernel launch failure:} In debug builds, the kernel may fail to launch but errors are not checked, while pre-initialized output buffers contain values matching CPU results.
    
    \item \textbf{Missing architecture guards:} Device code lacks \texttt{\_\_CUDA\_ARCH\_\_} preprocessor guards to provide fallback implementations for older compute capabilities.
\end{enumerate}

\subsubsection{Diagnostic Steps Performed}

\begin{itemize}
    \item Verified behavioral difference between debug (\texttt{-g -G}) and release (\texttt{-O3}) builds
    \item Observed missing kernel \texttt{printf} output despite ``successful'' execution
    \item Identified need for explicit error checking after kernel launch
    \item Recommended sentinel value initialization to verify kernel execution
\end{itemize}

\subsubsection{Recommended Fixes}

\begin{enumerate}
    \item \textbf{Add explicit error checking:}
    \begin{verbatim}
    kernel<<<grid, block>>>(args);
    cudaGetLastError(); cudaDeviceSynchronize();
    \end{verbatim}
    
    \item \textbf{Guard advanced instructions by compute capability:}
    \begin{verbatim}
    #if __CUDA_ARCH__ >= 800
      // Use cp.async (Ampere+)
    #else
      // Fallback: regular loads
    #endif
    \end{verbatim}
    
    \item \textbf{Target correct architecture:}
    \begin{verbatim}
    nvcc -gencode=arch=compute_70,code=sm_70 \
         -gencode=arch=compute_80,code=sm_80
    \end{verbatim}
    
    \item \textbf{Initialize output buffers with sentinel values} (e.g., \texttt{-999.0f}) to detect whether kernel modifies data
\end{enumerate}

\subsubsection{Workaround}

Compile with device debugging enabled: \texttt{nvcc -G -g -O0}. Note this disables optimizations and is unsuitable for performance measurements.


\subsection{cudaDMA Illegal Instruction Issue}

\subsubsection{Problem Description}

The \texttt{jacobi2D\_cudaDMA} kernel encounters an illegal instruction error when compiled with optimizations (\texttt{-O2}/\texttt{-O3}), but runs without errors when compiled with debug flags (\texttt{-g -G}). However, a critical observation reveals suspicious behavior even in the debug build:

\begin{itemize}
    \item \textbf{No illegal instruction error} in debug mode
    \item \textbf{No CPU-GPU mismatch} reported (results appear correct)
    \item \textbf{Print statements inside kernel do NOT appear} in console output
\end{itemize}

This suggests the kernel may not be executing at all, yet falsely reports success.

\subsubsection{Root Causes}

\begin{enumerate}
    \item \textbf{Architecture mismatch:} The \texttt{cudaDMAStrided} helper likely emits SM80+ instructions (e.g., \texttt{cp.async}) that are unsupported on older GPUs. Optimized builds inline these instructions without runtime guards.
    
    \item \textbf{Silent kernel launch failure:} In debug builds, the kernel may fail to launch but errors are not checked, while pre-initialized output buffers contain values matching CPU results.
    
    \item \textbf{Missing architecture guards:} Device code lacks \texttt{\_\_CUDA\_ARCH\_\_} preprocessor guards to provide fallback implementations for older compute capabilities.
\end{enumerate}

\subsubsection{Diagnostic Steps Performed}

\begin{itemize}
    \item Verified behavioral difference between debug (\texttt{-g -G}) and release (\texttt{-O3}) builds
    \item Observed missing kernel \texttt{printf} output despite ``successful'' execution
    \item Identified need for explicit error checking after kernel launch
    \item Recommended sentinel value initialization to verify kernel execution
\end{itemize}

\subsubsection{Recommended Fixes}

\begin{enumerate}
    \item \textbf{Add explicit error checking:}
    \begin{verbatim}
    kernel<<<grid, block>>>(args);
    cudaGetLastError(); cudaDeviceSynchronize();
    \end{verbatim}
    
    \item \textbf{Guard advanced instructions by compute capability:}
    \begin{verbatim}
    #if __CUDA_ARCH__ >= 800
      // Use cp.async (Ampere+)
    #else
      // Fallback: regular loads
    #endif
    \end{verbatim}
    
    \item \textbf{Target correct architecture:}
    \begin{verbatim}
    nvcc -gencode=arch=compute_70,code=sm_70 \
         -gencode=arch=compute_80,code=sm_80
    \end{verbatim}
    
    \item \textbf{Initialize output buffers with sentinel values} (e.g., \texttt{-999.0f}) to detect whether kernel modifies data
\end{enumerate}

\subsubsection{Workaround}

Compile with device debugging enabled: \texttt{nvcc -G -g -O0}. Note this disables optimizations and is unsuitable for performance measurements.


\subsection{cudaDMA Illegal Instruction Issue}

\subsubsection{Problem Description}

The \texttt{jacobi2D\_cudaDMA} kernel encounters an illegal instruction error when compiled with optimizations (\texttt{-O2}/\texttt{-O3}), but runs without errors when compiled with debug flags (\texttt{-g -G}). However, a critical observation reveals suspicious behavior even in the debug build:

\begin{itemize}
    \item \textbf{No illegal instruction error} in debug mode
    \item \textbf{No CPU-GPU mismatch} reported (results appear correct)
    \item \textbf{Print statements inside kernel do NOT appear} in console output
\end{itemize}

This suggests the kernel may not be executing at all, yet falsely reports success.

\subsubsection{Root Causes}

\begin{enumerate}
    \item \textbf{Architecture mismatch:} The \texttt{cudaDMAStrided} helper likely emits SM80+ instructions (e.g., \texttt{cp.async}) that are unsupported on older GPUs. Optimized builds inline these instructions without runtime guards.
    
    \item \textbf{Silent kernel launch failure:} In debug builds, the kernel may fail to launch but errors are not checked, while pre-initialized output buffers contain values matching CPU results.
    
    \item \textbf{Missing architecture guards:} Device code lacks \texttt{\_\_CUDA\_ARCH\_\_} preprocessor guards to provide fallback implementations for older compute capabilities.
\end{enumerate}

\subsubsection{Diagnostic Steps Performed}

\begin{itemize}
    \item Verified behavioral difference between debug (\texttt{-g -G}) and release (\texttt{-O3}) builds
    \item Observed missing kernel \texttt{printf} output despite ``successful'' execution
    \item Identified need for explicit error checking after kernel launch
    \item Recommended sentinel value initialization to verify kernel execution
\end{itemize}

\subsubsection{Recommended Fixes}

\begin{enumerate}
    \item \textbf{Add explicit error checking:}
    \begin{verbatim}
    kernel<<<grid, block>>>(args);
    cudaGetLastError(); cudaDeviceSynchronize();
    \end{verbatim}
    
    \item \textbf{Guard advanced instructions by compute capability:}
    \begin{verbatim}
    #if __CUDA_ARCH__ >= 800
      // Use cp.async (Ampere+)
    #else
      // Fallback: regular loads
    #endif
    \end{verbatim}
    
    \item \textbf{Target correct architecture:}
    \begin{verbatim}
    nvcc -gencode=arch=compute_70,code=sm_70 \
         -gencode=arch=compute_80,code=sm_80
    \end{verbatim}
    
    \item \textbf{Initialize output buffers with sentinel values} (e.g., \texttt{-999.0f}) to detect whether kernel modifies data
\end{enumerate}

\subsubsection{Workaround}

Compile with device debugging enabled: \texttt{nvcc -G -g -O0}. Note this disables optimizations and is unsuitable for performance measurements.


\subsection{cudaDMA Illegal Instruction Issue}

\subsubsection{Problem Description}

The \texttt{jacobi2D\_cudaDMA} kernel encounters an illegal instruction error when compiled with optimizations (\texttt{-O2}/\texttt{-O3}), but runs without errors when compiled with debug flags (\texttt{-g -G}). However, a critical observation reveals suspicious behavior even in the debug build:

\begin{itemize}
    \item \textbf{No illegal instruction error} in debug mode
    \item \textbf{No CPU-GPU mismatch} reported (results appear correct)
    \item \textbf{Print statements inside kernel do NOT appear} in console output
\end{itemize}

This suggests the kernel may not be executing at all, yet falsely reports success.

\subsubsection{Root Causes}

\begin{enumerate}
    \item \textbf{Architecture mismatch:} The \texttt{cudaDMAStrided} helper likely emits SM80+ instructions (e.g., \texttt{cp.async}) that are unsupported on older GPUs. Optimized builds inline these instructions without runtime guards.
    
    \item \textbf{Silent kernel launch failure:} In debug builds, the kernel may fail to launch but errors are not checked, while pre-initialized output buffers contain values matching CPU results.
    
    \item \textbf{Missing architecture guards:} Device code lacks \texttt{\_\_CUDA\_ARCH\_\_} preprocessor guards to provide fallback implementations for older compute capabilities.
\end{enumerate}

\subsubsection{Diagnostic Steps Performed}

\begin{itemize}
    \item Verified behavioral difference between debug (\texttt{-g -G}) and release (\texttt{-O3}) builds
    \item Observed missing kernel \texttt{printf} output despite ``successful'' execution
    \item Identified need for explicit error checking after kernel launch
    \item Recommended sentinel value initialization to verify kernel execution
\end{itemize}

\subsubsection{Recommended Fixes}

\begin{enumerate}
    \item \textbf{Add explicit error checking:}
    \begin{verbatim}
    kernel<<<grid, block>>>(args);
    cudaGetLastError(); cudaDeviceSynchronize();
    \end{verbatim}
    
    \item \textbf{Guard advanced instructions by compute capability:}
    \begin{verbatim}
    #if __CUDA_ARCH__ >= 800
      // Use cp.async (Ampere+)
    #else
      // Fallback: regular loads
    #endif
    \end{verbatim}
    
    \item \textbf{Target correct architecture:}
    \begin{verbatim}
    nvcc -gencode=arch=compute_70,code=sm_70 \
         -gencode=arch=compute_80,code=sm_80
    \end{verbatim}
    
    \item \textbf{Initialize output buffers with sentinel values} (e.g., \texttt{-999.0f}) to detect whether kernel modifies data
\end{enumerate}

\subsubsection{Workaround}

Compile with device debugging enabled: \texttt{nvcc -G -g -O0}. Note this disables optimizations and is unsuitable for performance measurements.
